% 原书物理页101--110；逐页图像自校，尚未独立复核.
\section{Galois 扩张与 Galois 对应}\label{sec:galois-2-7}
\subsection{定义}\label{sec:galois-2-7-1}
\begin{definition*}有限扩张 $K\supset k$ 若正规且可分，则称为\newterm{galois-extension}{Galois 扩张}{Galois extension}. 记 $\operatorname{Gal}(K/k):=\operatorname{Aut}(K/k)$ 为 $k$-代数 $K$ 的自同构群，称为 $K$ 在 $k$ 上的\newterm{galois-group}{Galois 群}{Galois group}.\end{definition*}
下面的定理汇集了 Galois 扩张的若干有用刻画.
\subsection{Galois 扩张的刻画}\label{sec:galois-2-7-2}
\begin{theorem*}设 $K\supset k$ 为有限扩张. 下列条件等价：
\begin{enumerate}[label=\roman*)]
\item $K$ 在 $k$ 上是 Galois 扩张.
\item 对每个 $\alpha\in K$，在 $K[X]$ 中有 $f_\alpha=\prod_{\beta\in\operatorname{Aut}(K/k)\cdot\alpha}(X-\beta)$.
\item $K$ 是一个可分多项式的分裂域.
\item $|\operatorname{Aut}_{k\text{-alg}}(K)|=[K:k]$.
\item $K^{\operatorname{Aut}(K/k)}=k$（$\operatorname{Aut}(K/k)$ 在 $K$ 上作用的不动点）.
\end{enumerate}\end{theorem*}
\begin{proof}
(i)$\Leftrightarrow$(ii). 按定义，$K$ 在 $k$ 上 Galois 当且仅当对每个 $\alpha\in K$，$f_\alpha$ 可分且在 $K[X]$ 中分裂. 此前已见，$\operatorname{Aut}(\bar k/k)$ 传递地置换 $f_\alpha$ 在代数闭包 $\bar k$ 中的根. 将 $K$ 嵌入 $\bar k$. 因 $\operatorname{Aut}(\bar k/k)$ 保持 $K$，$\operatorname{Aut}(K/k)$ 也传递地置换 $f_\alpha$ 在 $K$ 中的根，故 $\operatorname{Aut}(K/k)\alpha=\{f_\alpha\text{ 在 }K\text{ 中的根}\}$.
(i)$\Leftrightarrow$(iii). 已见可分多项式 $f$ 的分裂域是正规扩张（推论~\ref{sec:galois-2-3-5}），也是可分扩张（推论~\ref{sec:galois-2-5-4} 后的讨论）. 反之，若 $K$ 是 Galois 扩张，则因其可分而为单生成，并且包含生成元 $\alpha$ 的极小多项式 $f_\alpha$ 的所有根. 所以它是 $f_\alpha$ 的分裂域.
(i)$\Rightarrow$(iv). 设 $\bar k$ 是 $k$ 的代数闭包. 因 $K/k$ 可分，有 $|\operatorname{Hom}_{k\text{-alg}}(K,\bar k)|=[K:k]$. 固定嵌入 $\iota_1\in\operatorname{Hom}_{k\text{-alg}}(K,\bar k)$，考虑 $\operatorname{Aut}_{k\text{-alg}}(K)\to\operatorname{Hom}_{k\text{-alg}}(K,\bar k),\ \sigma\mapsto\iota_1\circ\sigma$. 因 $\iota_1$ 单射，该映射单射. 它也满射，因为任一其他嵌入 $\iota_2:K\hookrightarrow\bar k$ 与 $\iota_1$ 在 $\bar k$ 中有相同的像 $K'$，故复合 $\sigma:=\iota_1^{-1}\circ\iota_2$ 是 $K$ 的自同构，满足 $\iota_2=\iota_1\circ\sigma$. 所以 $|\operatorname{Aut}(K/k)|=[K:k]$.
(iv)$\Rightarrow$(v). 记 $k':=K^{\operatorname{Aut}(K/k)}$，显然它是 $K$ 中包含 $k$ 的子域. 所以 $\operatorname{Aut}(K/k')=\operatorname{Aut}(K/k)$. 取嵌入 $\iota:K\hookrightarrow\overline{k'}$ 至 $k'$ 的一个代数闭包. 映射 $\sigma\mapsto\iota\circ\sigma$ 将 $\operatorname{Aut}(K/k')$ 单射到 $\operatorname{Hom}_{k'\text{-alg}}(K,\overline{k'})$. 由命题~\ref{sec:galois-2-5-4}，$|\operatorname{Aut}(K/k')|\leq[K:k']$. 又 $[K:k']\leq[K:k]=|\operatorname{Aut}(K/k)|=|\operatorname{Aut}(K/k')|$. 所以 $[K:k']=[K:k]$，继而 $[k':k]=1$，最终 $k'=k$.
(v)$\Rightarrow$(ii). 取 $\alpha\in K$，置 $g_\alpha:=\prod_{\beta\in\operatorname{Aut}(K/k)\cdot\alpha}(X-\beta)\in K[X]$. 每个 $\beta\in\operatorname{Aut}(K/k)\alpha$ 都是 $f_\alpha$ 的根，故 $g_\alpha$ 整除 $f_\alpha$. 因 $f_\alpha$ 在 $k[X]$ 中不可约，只需证明 $g_\alpha\in k[X]$. 如常将 $G:=\operatorname{Aut}(K/k)$ 的作用延拓至 $K[X]$，即作用于多项式系数. 在假设(v)下，$f\in K[X]$ 属于 $k[X]$ 当且仅当它被 $G$ 固定. 对每个 $\sigma\in G$，
\begin{align*}
\sigma(g_\alpha)=\prod_{\beta\in G\alpha}(X-\sigma(\beta))=\prod_{\gamma\in\sigma G\alpha}(X-\gamma)=\prod_{\gamma\in G\alpha}(X-\gamma)=g_\alpha.
\end{align*}
所以 $g_\alpha$ 被 $G$ 固定，属于 $k[X]$.
\end{proof}
\subsection{例：有限域}\label{sec:galois-2-7-3}
\begin{example*}扩张 $\mathbb F_{p^r}\supset\mathbb F_p$ 是 Galois 扩张，因为它是可分多项式 $X^{p^r}-X$ 的分裂域. 令 $F$ 为 $\mathbb F_{p^r}$ 的 Frobenius 自同态，它是自同构，故属于 $\operatorname{Gal}(\mathbb F_{p^r}/\mathbb F_p)$. 显然 $F^r=\mathrm{id}$. 又对 $s<r$，此前已见不动点子域 $\mathbb F_{p^r}^{F^s}$ 是 $X^{p^s}-X$ 的根集，故基数小于 $p^r$. 所以 $F$ 的阶为 $r$，$\operatorname{Gal}(\mathbb F_{p^r}/\mathbb F_p)$ 是由 $F$ 生成的 $r$ 阶循环群.\end{example*}
\subsection{例：分圆扩张}\label{sec:galois-2-7-4}
域 $k$ 的 $n$ 次\newterm{galois-cyclotomic-extension}{分圆扩张}{cyclotomic extension}是多项式 $X^n-1$ 的分裂域 $k_n$，即由 $n$ 次单位根生成的扩张. 若 $k$ 的特征为 $p>0$，且 $n=p^kn'$、$(n',p)=1$，则 $X^n-1=(X^{n'}-1)^{p^k}$，故 $k_n=k_{n'}$. 因此不失一般性，假定 $(n,p)=1$. 因 $X^n-1$ 可分，$k_n\supset k$ 是 Galois 扩张. 求 $\operatorname{Gal}(k_n/k)$. 每个 $\sigma$ 显然保持 $k_n$ 中的 $n$ 次单位根子群 $\mu_n(k_n)$，且因该子群在 $k$ 上生成 $k_n$，$\sigma$ 完全由其在 $\mu_n(k_n)$ 上的作用确定. 又 $\sigma$ 通过群自同构作用于 $\mu_n(k_n)$，得到单射
\begin{align*}
\operatorname{Gal}(k_n/k)\hookrightarrow\operatorname{Aut}_{\mathrm{grp}}(\mu_n(k_n)).
\end{align*}
在定理~\ref{sec:galois-2-5-5} 的证明中已见，$\mu_n(k_n)$ 是 $n$ 阶循环群. 若 $\zeta_n$ 为其生成元（即本原 $n$ 次单位根），则 $\sigma(\zeta_n)$ 也是生成元，故为 $\zeta_n^{a_\sigma}$，其中 $a_\sigma\in\mathbb Z$ 模 $n$ 唯一确定，且 $(a_\sigma,n)=1$. 得到单射
\begin{align*}
\chi_k:\operatorname{Gal}(k_n/k)\hookrightarrow(\mathbb Z/n\mathbb Z)^\times,\qquad\sigma\mapsto(a_\sigma\bmod n).
\end{align*}
若无关于 $k$ 的进一步信息，就不能再多说什么. 下面给出几个例子.
\begin{itemize}
\item $k=\mathbb F_p$. 此时 $k_n$ 必形如 $\mathbb F_{p^r}=k_{p^r-1}$. 所以 $r$ 是满足 $n\mid p^r-1$ 的最小整数，即 $p$ 在 $(\mathbb Z/n\mathbb Z)^\times$ 中的阶. 已见 $\operatorname{Gal}(\mathbb F_{p^r}/\mathbb F_p)$ 是由 Frobenius $F$ 生成的 $r$ 阶循环群. 故 $\operatorname{Gal}(k_n/k)$ 也是如此，且由 Frobenius 和 $\chi$ 的定义，有 $\chi_{\mathbb F_p}(F)=(p\bmod n)$.
\item $k=\mathbb Q$. 于是 $\mathbb Q_n=\mathbb Q(e^{2i\pi/n})$. 以 $c$ 表示复共轭，它是 $\mathbb C$ 的自同构，必保持代数闭子域 $\overline{\mathbb Q}$ 及正规子扩张 $\mathbb Q_n$. 因为 $c$ 将 $e^{2i\pi/n}$ 映为 $e^{-2i\pi/n}$，有 $\chi_{\mathbb Q}(c)=(-1\bmod n)$. 实际上将证明：
\begin{theorem*}$\chi_{\mathbb Q}$ 是同构 $\operatorname{Gal}(\mathbb Q_n/\mathbb Q)\xrightarrow{\sim}(\mathbb Z/n\mathbb Z)^\times$.\end{theorem*}
由前面讨论，这等价于 $[\mathbb Q(e^{2i\pi/n}):\mathbb Q]=\varphi(n):=|(\mathbb Z/n\mathbb Z)^\times|$（Euler 函数）. 为此记
\begin{align*}
\Phi_n(X):=\prod_{\substack{0\leq i<n\\(a,n)=1}}(X-e^{2\pi ia/n})=\prod_{\zeta\text{ 的阶为 }n}(X-\zeta)\in\overline{\mathbb Q}[X].
\end{align*}
\footnote{译注：原书第一个乘积的界使用 $i$，互素条件及指数使用 $a$，指标混用，保留原式.}
第二个乘积遍历本原 $n$ 次单位根. 所以在 $\overline{\mathbb Q}[X]$ 中有分解 $X^n-1=\prod_{d\mid n}\Phi_d(X)$. 实际上，$\Phi_n(X)\in\mathbb Q_n[X]$ 在 $\operatorname{Gal}(\mathbb Q_n/\mathbb Q)$ 下不变，因为本原 $n$ 次单位根的任一共轭仍是本原 $n$ 次单位根. 由定理(v)，$\Phi_n(X)\in\mathbb Q[X]$. 所以 $\mathbb Q_n$ 也是 $\Phi_n$ 的分裂域，且因 $\deg\Phi_n=\varphi(n)$，只需证明：
\begin{lemma*}$\Phi_n$ 在 $\mathbb Q[X]$ 中不可约.\end{lemma*}
\begin{proof}设 $\Phi_n=fg$ 为 $\mathbb Q[X]$ 中的分解，$f,g$ 首一且 $\deg f>0$. 此分解对应于本原 $n$ 次单位根集合的一个划分. 要证明 $f=\Phi_n$，只需证明 $f$ 的根集在 $(\mathbb Z/n\mathbb Z)^\times$ 的作用下稳定，即对每个与 $n$ 互素的 $a$，取 $a$ 次幂后仍稳定. 当然只须对每个与 $n$ 互素的素数 $p$ 证明取 $p$ 次幂后稳定.
先证明 $\Phi_n\in\mathbb Z[X]$. 其系数属于 $\mathbb Q$ 的子环 $\mathbb Z[e^{2\pi i/n}]\cap\mathbb Q$. 而 $\mathbb Z[e^{2i\pi/n}]$ 作为 $\mathbb Z$-模由 $e^{2\pi im/n}$、$0\leq m<n$ 生成（参见~\ref{sec:galois-1-4-3} 的第二个推论）. 故子环 $\mathbb Z[e^{2\pi i/n}]\cap\mathbb Q$ 也是有限生成 $\mathbb Z$-模，从而等于 $\mathbb Z$.
再证明 $f,g\in\mathbb Z[X]$. 采用推论~\ref{sec:galois-1-7-2} 及其后注记的记号，只须证明对每个素数 $p$，$\nu_p(f),\nu_p(g)\geq0$. 因 $f,g$ 首一，有 $\nu_p(f),\nu_p(g)\leq0$；又 $\nu_p(f)+\nu_p(g)=0$，所以二者均为 $0$.
固定与 $n$ 互素的素数 $p$，以 $\overline\Phi_n,\bar f,\bar g$ 表示它们在 $\mathbb F_p[X]$ 中的像. $\overline\Phi_n$ 在 $\overline{\mathbb F}_p$ 中的根集是本原 $n$ 次单位根的集合，分解 $\overline\Phi_n=\bar f\bar g$ 仍对应于此集合的划分. 特别地，$\bar f$ 与 $\bar g^p$ 互素，故存在 $\bar u,\bar v\in\mathbb F_p[X]$ 使 $\bar u\bar f+\bar v\bar g^p=1$. 取提升 $u,v\in\mathbb Z[X]$，注意 $\overline{g(X^p)}=\bar g^p$，可知存在 $w\in\mathbb Z[X]$ 使
\begin{align*}uf(X)+vg(X^p)=1+pw.\end{align*}
取 $f$ 的根 $\zeta$. 有 $1+pw(\zeta)\ne0$，否则 $w(\zeta)=-1/p\in\mathbb Z[\zeta]\cap\mathbb Q=\mathbb Z$，矛盾. 所以 $\zeta^p$ 不可能是 $g$ 的根，只能是 $f$ 的根，如所欲证.\end{proof}
\end{itemize}
\subsection{例：根式扩张}\label{sec:galois-2-7-5}
设 $a\in k^\times$、$n\in\mathbb N$. 以 $\mu_n$ 记 $n$ 次单位根群，假设 $|\mu_n|=n$（即 $k$ 含有所有 $n$ 次单位根）. 考虑由满足 $\alpha^n=a$（即 $f_\alpha\mid X^n-a$）的元素 $\alpha$ 生成的 $k$ 的扩张 $K$. 则 $K$ 是 $X^n-a$ 在 $k$ 上的分裂域，因为每个其他根均形如 $\alpha\zeta$、$\zeta\in\mu_n$，因而属于 $K$. 由此还可得到 $\operatorname{Gal}(K/k)$ 的信息. 每个 $\sigma$ 完全由其在 $\alpha$ 上的作用确定，且 $\sigma(\alpha)=\alpha\zeta_\sigma$，其中 $\zeta_\sigma\in\mu_n$. 对另一 $\sigma'$，有 $(\sigma'\sigma)(\alpha)=\sigma'(\alpha\zeta_\sigma)=\sigma'(\alpha)\zeta_\sigma=\alpha\zeta_{\sigma'}\zeta_\sigma$，故映射 $\operatorname{Gal}(K/k)\to\mu_n,\ \sigma\mapsto\zeta_\sigma$ 是单射群同态. 因 $\mu_n$ 的每个子群均为某个 $m\mid n$ 的 $\mu_m$，得到同构 $\operatorname{Gal}(K/k)\xrightarrow{\sim}\mu_m$，其中 $m:=[K:k]$；特别地 Galois 群是循环群.
尝试由 $a$ 判断 $m$. 元素 $N\alpha=\prod_\sigma\sigma(\alpha)$ 属于 $K^{\operatorname{Gal}(K/k)}=k$. 有 $N\alpha=\alpha^m\prod_\sigma\zeta_\sigma$，所以 $\alpha^m\in k$，继而 $a=(\alpha^m)^{n/m}\in(k^\times)^{n/m}$. 因此 $m$ 是 $a$ 在商群 $k^\times/(k^\times)^n$ 中的阶；特别地，$X^n-a$ 在 $k[X]$ 中不可约，当且仅当 $a$ 在此商群中的阶为 $n$. 综上已证明下面定理的第一部分.
\begin{theorem*}设 $k$ 为满足 $|\mu_n(k)|=n$ 的域. 对每个 $a\in k$，由 $a$ 的一个 $n$ 次根生成的扩张 $k[\sqrt[n]a]$ 是次数 $m$ 的 Galois 扩张，其中 $m$ 是 $a$ 在 $k^\times/(k^\times)^n$ 中的阶，且有同构
\begin{align*}
\operatorname{Gal}(k[\sqrt[n]a]/k)\xrightarrow{\sim}\mu_m,\qquad\sigma\mapsto\frac{\sigma(\sqrt[n]a)}{\sqrt[n]a}.
\end{align*}
反之，Galois 群为 $n$ 阶循环群的每个扩张 $K\supset k$ 均形如 $k[\sqrt[n]a]$，其中 $a\in k$.\end{theorem*}
\noindent\emph{译注：}原书定理量词写 $a\in k$，但其阶及上述分式要求 $a\ne0$；保留原陈述.
只须证明逆命题. 利用事实：$\operatorname{Gal}(K/k)=\operatorname{Hom}_{k\text{-alg}}(K,K)$ 在 $K$-向量空间 $\operatorname{Hom}_{k\text{-ev}}(K,K)$ 中线性无关. 为说明这一点，只须看出 $\operatorname{Hom}_{k\text{-alg}}(K,\bar k)=\operatorname{Hom}_{\bar k\text{-alg}}(\bar k\otimes_kK,\bar k)$ 在 $\bar k$-向量空间 $\operatorname{Hom}_{k\text{-ev}}(K,\bar k)=\operatorname{Hom}_{\bar k\text{-ev}}(\bar k\otimes_kK,\bar k)$ 中线性无关；引理~\ref{sec:galois-2-5-4} 的满射性说明，对 $\bar k$ 上任意有限维代数 $A$，$\operatorname{Hom}_{\bar k\text{-alg}}(A,\bar k)$ 在 $\operatorname{Hom}_{\bar k\text{-ev}}(A,\bar k)$ 中线性无关.
取 $\operatorname{Gal}(K/k)$ 的生成元 $\sigma$，于是群为 $\{1,\sigma,\ldots,\sigma^{n-1}\}$；取本原 $n$ 次单位根 $\zeta\in K$. $\sigma^i$ 的线性无关性保证存在 $x\in K$ 使 $\alpha:=x+\zeta^{-1}\sigma(x)+\cdots+\zeta^{1-n}\sigma^{n-1}(x)\ne0$. 则 $\sigma(\alpha)=\zeta\alpha$，故 $\sigma^i(\alpha)=\zeta^i\alpha$、$0\leq i<n$ 两两不同，且
\begin{align*}f_\alpha=\prod_{i=0}^{n-1}(X-\sigma^i(\alpha))=\prod_{i=0}^{n-1}(X-\zeta^i\alpha)=X^n-\alpha^n.\end{align*}
特别地，$\alpha^n\in k^\times$，且因 $f_\alpha$ 次数为 $n$，有 $K=k[\alpha]$，如所欲证.
下面的直接推论以后会用到.
\begin{corollary*}设 $k$ 为满足 $|\mu_n(k)|=n$ 的域，$K\supset k$ 由满足 $\alpha_i^n\in k$ 的元素 $\alpha_1,\ldots,\alpha_r$ 生成. 则 $K\supset k$ 为 Galois 扩张，其 Galois 群是 Abel 群.\end{corollary*}
\begin{proof}此扩张正规，因为它含有生成元 $\alpha_i$ 的所有共轭（形如 $\alpha_i\zeta_n^j$）. 故它是多项式 $(X^n-\alpha_1)\cdots(X^n-\alpha_r)$ 的分裂域.\footnote{译注：原书此式常数项写为 $\alpha_i$，与条件 $\alpha_i^n\in k$ 不合；保留原式.} 此扩张由可分元素生成，因而可分，所以是 Galois 扩张. 考虑映射
\begin{align*}\operatorname{Gal}(K/k)\to\prod_{i=1}^r\operatorname{Gal}(k(\alpha_i)/k),\quad\sigma\mapsto(\sigma|_{k(\alpha_1)},\ldots,\sigma|_{k(\alpha_r)}).\end{align*}
因各 $k(\alpha_i)/k$ 均 Galois，此映射良定义；因 $\alpha_i$ 生成 $K$，它单射，且为群同态. 所以 $\operatorname{Gal}(K/k)$ 是循环群之积的子群，从而是 Abel 群.\end{proof}
\subsection{逆问题}\label{sec:galois-2-7-6}
上述例子的 Galois 群都是 Abel 群. 下述结论是定理~\ref{sec:galois-2-7-2}(v) 的直接而有用的推论，它说明每个有限群都是一个 Galois 群.
\begin{corollary*}设 $G$ 是域 $K$ 的一个有限自同构群. 则 $K^G$ 是域，且 $K\supset K^G$ 为 Galois 扩张，满足 $\operatorname{Gal}(K/K^G)=G$.\end{corollary*}
\begin{proof}显然 $G\subset\operatorname{Aut}(K/K^G)$. 故 $K^G\subset K^{\operatorname{Aut}(K/K^G)}\subset K^G$. 只须应用定理(v).\end{proof}
\noindent\emph{译注：}原书此略证未先证明 $K/K^G$ 有限，也未由所写不动域等式单独论证自同构群恰为 $G$；此处保留略证.
\begin{example*}置换群 $\mathfrak S_n$ 通过置换不定元作用于 $K_n:=k(X_1,\ldots,X_n)$. 每个有限群 $G$ 都嵌入某个 $\mathfrak S_n$（例如取 $n=|G|$，令 $G$ 通过左平移作用于自身）. 则 $K_n\supset K_n^G$ 是群为 $G$ 的 Galois 扩张. 因而每个有限群都是 Galois 群. 仍未解决的 Galois 逆问题问：哪些有限群 $G$ 能成为某个 Galois 扩张 $K\supset\mathbb Q$ 的 Galois 群（人们认为所有有限群都可以）.\end{example*}

\subsection{Galois 对应}\label{sec:galois-2-7-7}
我们将建立 Galois 扩张的子扩张与其 Galois 群的子群之间的一个非凡双射，即\newterm{galois-correspondence}{Galois 对应}{Galois correspondence}. 先证明如下结果.
\begin{proposition*}设 $K\supset k$ 是 Galois 扩张，$k\subset K'\subset K$ 为子扩张. 则：
\begin{enumerate}[label=\roman*)]
\item $K$ 在 $K'$ 上 Galois，且 $K'=K^{\operatorname{Gal}(K/K')}$.
\item 对每个 $\sigma\in\operatorname{Gal}(K/k)$，有 $\operatorname{Gal}(K/\sigma(K'))=\sigma\operatorname{Gal}(K/K')\sigma^{-1}$.
\item $K'$ 在 $k$ 上 Galois，当且仅当 $\operatorname{Gal}(K/K')$ 是 $\operatorname{Gal}(K/k)$ 的正规子群. 此时 $\sigma\mapsto\sigma|_{K'}$ 诱导同构
\begin{align*}\operatorname{Gal}(K/k)/\operatorname{Gal}(K/K')\xrightarrow{\sim}\operatorname{Gal}(K'/k).\end{align*}
\end{enumerate}\end{proposition*}
\begin{proof}(i) 若 $K$ 是可分多项式 $f\in k[X]$ 在 $k$ 上的分裂域，则也是同一多项式在 $K'$ 上的分裂域，而 $f$ 仍可分. 故 $K$ 在 $K'$ 上 Galois，且 $K'=K^{\operatorname{Gal}(K/K')}$ 由定理(v)得出.
(ii) 对 $\tau\in\operatorname{Gal}(K/k)$，有
\begin{align*}
\tau\in\operatorname{Gal}(K/\sigma(K'))&\Longleftrightarrow(\forall\alpha\in K',\ \tau(\sigma(\alpha'))=\sigma(\alpha'))\\
&\Longleftrightarrow(\forall\alpha\in K',\ \sigma^{-1}\tau\sigma(\alpha)=\alpha)\\
&\Longleftrightarrow\sigma^{-1}\tau\sigma\in\operatorname{Gal}(K/K').
\end{align*}
\footnote{译注：原书第一行量词为 $\alpha$、被作用元为 $\alpha'$，指标不一致，保留原式.}
(iii) 若 $\operatorname{Gal}(K/K')$ 在 $\operatorname{Gal}(K/k)$ 中正规，则由(ii)，
\begin{align*}
\forall\sigma\in\operatorname{Gal}(K/k),\quad\sigma(K')=K^{\operatorname{Gal}(K/\sigma(K'))}=K^{\sigma\operatorname{Gal}(K/K')\sigma^{-1}}=K^{\operatorname{Gal}(K/K')}=K'.
\end{align*}
所以 $K'$ 在 $k$ 上正规. 它又可分，故是 Galois 扩张. 反之，设 $K'$ 在 $k$ 上 Galois. 则 $K/k$ 的每个自同构都保持 $K'$，从而诱导 $K'/k$ 的自同构. 限制自同构得到群同态 $\operatorname{Gal}(K/k)\to\operatorname{Gal}(K'/k),\ \sigma\mapsto\sigma|_{K'}$. 其核由在 $K'$ 上为恒等的 $K/k$ 自同构构成，即 $\operatorname{Gal}(K/K')$，因而后者正规. 为见此同态满射，可将 $K$ 嵌入代数闭包 $\bar k$，并回忆限制 $\operatorname{Gal}(\bar k/k)\to\operatorname{Gal}(K'/k)$ 满射. 也可注意到像的基数为 $[K:k][K:K']^{-1}=[K':k]$.
\end{proof}
固定有限 Galois 扩张 $K\supset k$. 以 $\mathcal{SE}(K)$ 表示包含于 $K$ 的子扩张 $K'\supset k$ 的集合，按包含关系排序. 又记 $G:=\operatorname{Gal}(K/k)$，以 $\mathcal{SG}(G)$ 表示 $G$ 的子群集合，也按包含关系排序. 有两个显然反序的映射
\begin{align*}
\mathcal{SE}(K)&\longrightarrow\mathcal{SG}(G),& K'&\longmapsto\operatorname{Gal}(K/K'),\\
\mathcal{SG}(G)&\longrightarrow\mathcal{SE}(K),& G'&\longmapsto K^{G'}.
\end{align*}
\begin{theorem*}这两个映射互为逆双射，并将 Galois 子扩张与正规子群对应起来.\end{theorem*}
\begin{proof}由前述命题和推论即得.\end{proof}

\begin{example*}$X^3-2$ 在 $\mathbb Q$ 上的分裂域 $K$ 是 $\mathbb Q$ 的 Galois 扩张. 已见 $K=\mathbb Q(j,\sqrt[3]2)$ 在 $\mathbb Q$ 上次数为 $6$，因为它在 $\mathbb Q(\sqrt[3]2)$ 上次数为 $2$，后者在 $\mathbb Q$ 上次数为 $3$. 故 $\operatorname{Gal}(K/\mathbb Q)$ 是 $6$ 阶群. 它嵌入 $\{\sqrt[3]2,j\sqrt[3]2,j^2\sqrt[3]2\}$ 的置换群 $\mathfrak S_3$，而该群也有 $6$ 个元素，故 $\operatorname{Gal}(K/\mathbb Q)\simeq\mathfrak S_3$.
再给出一种可推广的描述. $\operatorname{Gal}(K/\mathbb Q)$ 含有 $3$ 阶子群 $\operatorname{Gal}(K/\mathbb Q(j))\simeq\mathbb Z/3\mathbb Z$ 和 $2$ 阶子群 $\operatorname{Gal}(K/\mathbb Q(\sqrt[3]2))\simeq\mathbb Z/2\mathbb Z$. 前者正规，因为 $\mathbb Q(j)/\mathbb Q$ Galois；后者则不正规. 因此有半直积
\begin{align*}
\operatorname{Gal}(K/\mathbb Q)=\operatorname{Gal}(K/\mathbb Q(\sqrt[3]2))\ltimes\operatorname{Gal}(K/\mathbb Q(j))\simeq\mathbb Z/2\mathbb Z\ltimes\mathbb Z/3\mathbb Z.
\end{align*}
由此得到中间扩张的结构：恰有三个 $3$ 次子扩张 $\mathbb Q(\sqrt[3]2)$、$\mathbb Q(j\sqrt[3]2)$、$\mathbb Q(j^2\sqrt[3]2)$，对应三个 $2$ 阶子群；还有一个 $2$ 次子扩张 $\mathbb Q(j)$，它是 Galois 扩张.
现在还能证明 $\alpha=j+\sqrt[3]2$ 是 $K$ 在 $\mathbb Q$ 上的生成元. 令 $\sigma$ 为 $\operatorname{Gal}(K/\mathbb Q(\sqrt[3]2))$ 的生成元，$\tau$ 为将 $\sqrt[3]2$ 映为 $j\sqrt[3]2$ 的 $\operatorname{Gal}(K/\mathbb Q(j))$ 的生成元. 计算得 $j+\sqrt[3]2$ 的共轭集合
\begin{align*}
\{\alpha,\sigma(\alpha),\tau(\alpha),\tau^2(\alpha),\sigma\tau(\alpha),\sigma\tau^2(\alpha)\}
\\=\bigl\{j+\sqrt[3]2,j^2+\sqrt[3]2,j+j\sqrt[3]2,j+j^2\sqrt[3]2,j^2+j^2\sqrt[3]2,j^2+j\sqrt[3]2\bigr\}
\end{align*}
基数为 $6$，故 $\deg f_\alpha=6$，$\alpha$ 在 $\mathbb Q$ 上生成 $K$.\end{example*}
下面推广这个例子.
\subsection{例：$\mathbb Q$ 上的 Kummer 扩张}\label{sec:galois-2-7-8}
固定 $n\in\mathbb N$、$n>1$，置 $\zeta_n:=e^{2\pi i/n}$. 设 $a\in\mathbb Q^\times$ 在商群 $\mathbb Q[\zeta_n]^\times/(\mathbb Q[\zeta_n]^\times)^n$ 中的像的阶为 $n$. 前面已见，此时 $X^n-a$ 在 $\mathbb Q[\zeta_n][X]$ 中不可约，更在 $\mathbb Q[X]$ 中不可约. 取 $a$ 在 $\mathbb C$ 中的一个 $n$ 次根 $\sqrt[n]a$. 域 $\mathbb Q[\sqrt[n]a]$ 是 $X^n-a$ 的根域，但只要 $n>2$，就不是 Galois 扩张. 而域 $\mathbb Q[\zeta_n,\sqrt[n]a]$ 是其分裂域，故在 $\mathbb Q$ 上 Galois. 它在 $\mathbb Q[\zeta_n]$ 上次数为 $n$，故在 $\mathbb Q$ 上次数为 $n\varphi(n)$. 群 $G:=\operatorname{Gal}(\mathbb Q[\zeta_n,\sqrt[n]a]/\mathbb Q)$ 含有两个值得注意的子群
\begin{align*}
H_1:=\operatorname{Gal}(\mathbb Q[\zeta_n,\sqrt[n]a]/\mathbb Q[\zeta_n]),\qquad H_a:=\operatorname{Gal}(\mathbb Q[\zeta_n,\sqrt[n]a]/\mathbb Q[\sqrt[n]a]).
\end{align*}
若 $n=2$，则 $H_1=\{\mathrm{id}\}$、$H_a=G$.\footnote{译注：原书如此，按前面的定义应当是 $H_1=G$、$H_a=\{\mathrm{id}\}$，这里保留原陈述.} 设 $n>2$，计算这些群. 在~\ref{sec:galois-2-7-5} 已构造同构
\begin{align*}H_1\xrightarrow{\sim}\mu_n,\qquad\sigma\mapsto\zeta_\sigma:=\frac{\sigma(\sqrt[n]a)}{\sqrt[n]a}.\end{align*}
又由 $[\mathbb Q[\zeta_n,\sqrt[n]a]:\mathbb Q[\sqrt[n]a]]\,[\mathbb Q[\sqrt[n]a]:\mathbb Q]=[\mathbb Q[\zeta_n,\sqrt[n]a]:\mathbb Q]=n\varphi(n)$ 得 $[\mathbb Q[\zeta_n,\sqrt[n]a]:\mathbb Q[\sqrt[n]a]]=\varphi(n)$，故~\ref{sec:galois-2-7-4} 中的同态
\begin{align*}\chi_{\mathbb Q[\sqrt[n]a]}:H_a\xrightarrow{\sim}(\mathbb Z/n\mathbb Z)^\times,\qquad\tau\mapsto a_\tau\end{align*}
是同构（顺便说明 $\Phi_n$ 在 $\mathbb Q[\sqrt[n]a][X]$ 中仍不可约）. 因 $\mathbb Q[\zeta_n]/\mathbb Q$ Galois，$H_1$ 在 $G$ 中正规，特别被 $H_a$ 正规化. 又 $H_a\cap H_1=\{\mathrm{id}\}$，因为交中的 $\sigma$ 固定 $\sqrt[n]a$ 和 $\zeta_n$，从而固定 $\mathbb Q[\zeta_n,\sqrt[n]a]$. 因 $|G|=|H_a||H_1|$，故 $G=H_1\rtimes H_a$ 是 $H_1$ 与 $H_a$ 的半直积. 明确写出 $H_a$ 在 $H_1$ 上的共轭作用. 取 $\sigma\in H_1$、$\tau\in H_a$，有
\begin{align*}(\tau\sigma\tau^{-1})(a)=(\tau\sigma)(a)=\tau(a\zeta_\sigma)=a\tau(\zeta_\sigma)=a\zeta_\sigma^{a_\tau}.\end{align*}
\footnote{译注：原书在这一行将被作用元素印为 $a$；$a\in\mathbb Q$ 被全部自同构固定，此处应涉及 $\sqrt[n]a$，保留原式.}
因此 $H_a$ 在 $H_1$ 上的共轭作用对应于 $(\mathbb Z/n\mathbb Z)^\times$ 在 $\mu_n$ 上的自然作用 $(a,\zeta)\mapsto\zeta^a$. 得到
\begin{align*}\operatorname{Gal}(\mathbb Q[\zeta_n,\sqrt[n]a]/\mathbb Q)\xrightarrow{\sim}\mu_n\rtimes(\mathbb Z/n\mathbb Z)^\times,\quad\sigma\mapsto\left(\frac{\sigma(\sqrt[n]a)}{\sqrt[n]a},\sigma|_{\mu_n}\right).\end{align*}
\begin{remark*}这里假设 $X^n-a$ 在 $\mathbb Q[\zeta_n][X]$ 中不可约，一般比在 $\mathbb Q[X]$ 中不可约更强. 但若 $n$ 为素数，或更一般地 $n$ 与 $\varphi(n)$ 互素，则二者等价. 因为此时必有 $\mathbb Q[\zeta_n]\cap\mathbb Q[\sqrt[n]a]=\mathbb Q$，可以应用下面的命题.\end{remark*}
\subsection{线性无交的子扩张与半直积}\label{sec:galois-2-7-9}
设 $K\supset k$ 为可分代数扩张，$K_1,K_2$ 是在 $k$ 上有限的子扩张.
\begin{proposition*}记 $K_{12}$ 为 $K$ 中由 $K_1,K_2$ 生成的子域. 下列四个性质等价：
\begin{enumerate}[label=\roman*)]
\item $[K_{12}:k]=[K_1:k][K_2:k]$.
\item $[K_{12}:K_1]=[K_2:k]$.
\item 典范同态 $K_1\otimes_kK_2\to K_{12},\ x_1\otimes x_2\mapsto x_1x_2$ 是同构.
\item 对每个 $\alpha\in K_2$，$\alpha$ 在 $k$ 上的极小多项式 $f_\alpha\in k[X]$ 在 $K_1[X]$ 中仍不可约.
\end{enumerate}
若 $K_1$ 或 $K_2$ 在 $k$ 上正规，则这些性质也等价于
\begin{enumerate}[label=\roman*),start=5]\item $K_1\cap K_2=k$.\end{enumerate}
\end{proposition*}
\begin{proof}(i)$\Leftrightarrow$(ii) 由 $[K_{12}:k]=[K_{12}:K_1][K_1:k]$ 得出. (ii)$\Leftrightarrow$(iii) 由(iii)中同态的满射性（按``生成的域''的定义）及标量扩张（此处从 $k$ 到 $K_1$）保持维数得出.
(iii)$\Rightarrow$(iv). (iii)中的同构在 $K_1\otimes_k k[\alpha]$ 上诱导到其像的同构，而其像恰为 $K_1[\alpha]$，说明 $\alpha$ 在 $K_1$ 和 $k$ 上的次数相同.
(iv)$\Rightarrow$(ii). 取 $\alpha$ 使 $k[\alpha]=K_2$，记其在 $k$ 上的极小多项式为 $f_\alpha$. 则 $[K_2:k]=\deg f_\alpha$. 又 $K_{12}=K_1[\alpha]$，且(iv)说明 $f_\alpha$ 也是 $\alpha$ 在 $K_1$ 上的极小多项式. 故 $[K_{12}:K_1]=\deg f_\alpha=[K_2:k]$.
(iv)$\Rightarrow$(v) 无须额外假设. 若 $\alpha\in K_1\cap K_2$，其在 $K_1$ 上的极小多项式为 $X-\alpha$. 由(iv)，此多项式属于 $k[X]$，所以 $\alpha\in k$.
(v)$\Rightarrow$(iv). 先注意，(iv)与(i)的等价性表明(iv)在交换 $K_1,K_2$ 后保持不变，这一点事先并不显然. 为确定起见，设 $K_2/k$ 正规. 则对 $\alpha\in K_2$，其在 $k$ 上的极小多项式 $f_\alpha$ 在 $K_2[X]$ 中分裂. 令 $g_\alpha\in K_1[X]$ 为 $\alpha$ 在 $K_1$ 上的极小多项式. 则 $g_\alpha\mid f_\alpha$，且 $g_\alpha\in K_2[X]$，因为其系数是 $g_\alpha$ 在 $K_2$ 中的根的多项式. 由此 $g_\alpha\in(K_1\cap K_2)[X]=k[X]$，故 $g_\alpha=f_\alpha$.\end{proof}
\begin{remark*}要使(v)蕴含其他性质，额外假设是必要的. 例如 $\mathbb Q(\sqrt[3]2)\cap\mathbb Q(j\sqrt[3]2)=\mathbb Q$，但 $\mathbb Q(\sqrt[3]2,j\sqrt[3]2)=\mathbb Q(\sqrt[3]2,j)$ 次数为 $6$ 而非 $9$.\end{remark*}
\begin{corollary*}在上述命题的情形下，设 $K_{12}$ 和 $K_1$ 在 $k$ 上 Galois. 则 $\operatorname{Gal}(K_{12}/k)$ 是其正规子群 $\operatorname{Gal}(K_{12}/K_1)$ 与其子群 $\operatorname{Gal}(K_{12}/K_2)$ 的半直积. 更确切地，$(\sigma,\tau)\mapsto\sigma\tau$ 是同构
\begin{align*}\operatorname{Gal}(K_{12}/K_1)\rtimes\operatorname{Gal}(K_{12}/K_2)\xrightarrow{\sim}\operatorname{Gal}(K_{12}/k),\end{align*}
其中半直积相对于共轭作用取定.\end{corollary*}
\noindent\emph{译注：}原书只写``在命题的情形下''，未明确假设命题中的等价性质成立；若不作此假设，结论一般不成立.
\begin{proof}$\operatorname{Gal}(K_{12}/K_2)$ 正规，因为 $K_2$ 是 Galois 扩张.\footnote{译注：原书此句与陈述中的 $K_1$ 写反，应涉及 $\operatorname{Gal}(K_{12}/K_1)$ 和 $K_1$；此处保留.} 交 $\operatorname{Gal}(K_{12}/K_2)\cap\operatorname{Gal}(K_{12}/K_1)$ 是固定 $K_1,K_2$，从而固定它们生成的域 $K_{12}$ 的自同构群. 所以交为 $\{\mathrm{id}\}$. 因而陈述中的映射单射. 因两集合基数相同，它双射. 最后，公式 $(\sigma\tau)(\sigma'\tau')=(\sigma\cdot\tau\sigma'\tau^{-1})(\tau\tau')$ 说明它是群同态.\end{proof}
\subsection{扩张的正规闭包与 Galois 闭包}\label{sec:galois-2-7-10}
多项式分裂域的概念在域扩张中也有类似物，即正规闭包.
\begin{definition*}设 $K\supset k$ 为代数扩张. 若扩张 $\widetilde K\supset k$ 正规，且由所有 $k$-嵌入 $\iota:K\hookrightarrow\widetilde K$ 的像 $\iota(K)$ 生成，则称其为 $K$ 的\newterm{galois-normal-closure}{正规闭包}{normal closure}.\end{definition*}
当 $K\supset k$ 为有限可分扩张时，也称 $\widetilde K\supset k$ 为 $K\supset k$ 的\newterm{galois-galois-closure}{Galois 闭包}{Galois closure}. 因为此时 $\widetilde K$ 在 $k$ 上也有限可分.
\begin{example*}不可约可分多项式 $f\in k[X]$ 的分裂域是其根域的一个 Galois 闭包.\end{example*}
\begin{proposition*}每个扩张都有正规闭包，在同构意义下唯一.\end{proposition*}
\begin{proof}取代数闭包 $\bar k$，以 $\widetilde K$ 记 $\bar k$ 中由所有像 $\iota(K)$、$\iota\in\operatorname{Hom}_{k\text{-alg}}(K,\bar k)$ 生成的子域. 它显然是 $K$ 的正规闭包. 若 $\widetilde K'$ 是另一个正规闭包，可将其嵌入 $\bar k$，其像必为 $\widetilde K$，因而得到同构 $\widetilde K'\xrightarrow{\sim}\widetilde K$.\end{proof}
另一种说法是：若 $K$ 已嵌入代数闭包 $\bar k$，它在 $K$ 中的正规闭包是由所有 $\sigma(K)$、$\sigma\in\operatorname{Aut}(\bar k/k)$ 生成的域.\footnote{译注：原书此处印为``在 $K$ 中''，应与前文所指代数闭包 $\bar k$ 区分；保留原说法.}
\begin{example*}若 $K=k(\alpha_1,\ldots,\alpha_n)$，其中 $\alpha_i\in\bar k$，则 $\widetilde K=K(\{\alpha_i^{(j)}\}_{i=1,\ldots,n;\,j=1,\ldots,r_j})$，其中 $\alpha_i^{(j)}$、$j=1,\ldots,r_i$ 表示 $\alpha_i$ 在 $\bar k$ 中的共轭. 换言之，$K$ 是多项式 $f_{\alpha_1}\cdots f_{\alpha_n}$ 的分裂域.\footnote{译注：原书生成集合的上界印为 $r_j$，随后说明用 $r_i$；末句主语印为 $K$ 而非 $\widetilde K$. 均保留原文.}\end{example*}
\begin{example*}$\mathbb Q(\sqrt[3]2,\sqrt[5]3)$ 在 $\overline{\mathbb Q}$ 中的 Galois 闭包是 $\mathbb Q(\sqrt[3]2,\sqrt[5]3,e^{2i\pi/15})$.\end{example*}



